\documentclass[10pt,a4paper]{amsart}
\usepackage{amsmath,amssymb,amsthm,graphicx}
% Retained preamble for the Stacks Project chapter snapshot in this directory.
% The chapter's own preamble.tex is a copy of the shared Stacks Project preamble; the block
% below is sliced MECHANICALLY, line for line, from it. Nothing was added, removed or
% redefined; the wrapper class, the xr-hyper/hyperref wiring and the externaldocument list
% are omitted because the assembler supplies the document class.
%
% lines taken from preamble.tex: 3, 5, 9, 11, 12, 156, 157, 158, ... (51 lines total)
%
% package coverage restored from the chapter's own preamble (2026-10-13): xy, xy-2cell

\usepackage[all]{xy}
\xyoption{2cell}
\UseAllTwocells

\newenvironment{reference}{\comment}{\endcomment}
\newenvironment{slogan}{\comment}{\endcomment}
\newenvironment{history}{\comment}{\endcomment}
\theoremstyle{plain}
\newtheorem{theorem}[subsection]{Theorem}
\newtheorem{proposition}[subsection]{Proposition}
\newtheorem{lemma}[subsection]{Lemma}
\theoremstyle{definition}
\newtheorem{definition}[subsection]{Definition}
\newtheorem{example}[subsection]{Example}
\newtheorem{exercise}[subsection]{Exercise}
\newtheorem{situation}[subsection]{Situation}
\theoremstyle{remark}
\newtheorem{remark}[subsection]{Remark}
\newtheorem{remarks}[subsection]{Remarks}
\numberwithin{equation}{subsection}
\def\lim{\mathop{\mathrm{lim}}\nolimits}
\def\colim{\mathop{\mathrm{colim}}\nolimits}
\def\Spec{\mathop{\mathrm{Spec}}}
\def\Hom{\mathop{\mathrm{Hom}}\nolimits}
\def\Ext{\mathop{\mathrm{Ext}}\nolimits}
\def\SheafHom{\mathop{\mathcal{H}\!\mathit{om}}\nolimits}
\def\SheafExt{\mathop{\mathcal{E}\!\mathit{xt}}\nolimits}
\def\Sch{\mathit{Sch}}
\def\Mor{\mathop{\mathrm{Mor}}\nolimits}
\def\Ob{\mathop{\mathrm{Ob}}\nolimits}
\def\Sh{\mathop{\mathit{Sh}}\nolimits}
\def\NL{\mathop{N\!L}\nolimits}
\def\CH{\mathop{\mathrm{CH}}\nolimits}
\def\proetale{{pro\text{-}\acute{e}tale}}
\def\etale{{\acute{e}tale}}
\def\QCoh{\mathit{QCoh}}
\def\Ker{\mathop{\mathrm{Ker}}}
\def\Im{\mathop{\mathrm{Im}}}
\def\Coker{\mathop{\mathrm{Coker}}}
\def\Coim{\mathop{\mathrm{Coim}}}
\DeclareMathSymbol{\boxtimes}{\mathbin}{AMSa}{"02}
\def\QCohstack{\mathcal{QC}\!\mathit{oh}}
\def\Cohstack{\mathcal{C}\!\mathit{oh}}
\def\Spacesstack{\mathcal{S}\!\mathit{paces}}
\def\Quotfunctor{\mathrm{Quot}}
\def\Hilbfunctor{\mathrm{Hilb}}
\def\Curvesstack{\mathcal{C}\!\mathit{urves}}
\def\Polarizedstack{\mathcal{P}\!\mathit{olarized}}
\def\Complexesstack{\mathcal{C}\!\mathit{omplexes}}
\def\Pic{\mathop{\mathrm{Pic}}\nolimits}
\def\Picardstack{\mathcal{P}\!\mathit{ic}}
\def\Picardfunctor{\mathrm{Pic}}
\def\Deformationcategory{\mathcal{D}\!\mathit{ef}}
\usepackage{cleveref}
\crefname{originallemma}{original lemma}{original lemmas}
\Crefname{originallemma}{Original lemma}{Original lemmas}
\makeatletter
\newlabel{algebra-section-netherlander}{{00S0}{source}{}{sourcepointer-1}{}}
\newlabel{algebra-section-netherlander@cref}{{[originallemma][0][]00S0}{[1][0][]source}{}{}{}}
\makeatother
\crefname{originallemma}{original lemma}{original lemmas}
\Crefname{originallemma}{Original lemma}{Original lemmas}
\makeatletter
\newlabel{equation-to-solve}{{08S4}{source}{}{sourcepointer-2}{}}
\newlabel{equation-to-solve@cref}{{[originallemma][0][]08S4}{[1][0][]source}{}{}{}}
\makeatother
\title{Stacks, Deformation Theory: lemma-huge-diagram}
\author{The Stacks Project Authors}
\date{stacks-project:defos.tex:local-snapshot-20261012; GNU FDL 1.2 or later}
\begin{document}
\maketitle
\noindent\textit{Editorial note.} Statement span is lines 274-297 of sources/stacks-defos/pinned-original-source.tex; proof span is lines 299-373, the complete author proof (73 body lines).\par\medskip
\noindent\textit{Editorial note.} Navigation stubs: 2 cited-but-not-included labels are supplied as non-load-bearing pointers with their pinned tags where the snapshot tag index has them.\par\medskip
\noindent\textit{Editorial note.} extra\_packages is empty on purpose: the retained preamble already loads xy with options and a second declaration clashes in the assembler wrapper.\par\medskip
\noindent\textit{Editorial note.} Text is retained verbatim, including the author's own typos and placeholders.\par\medskip
\noindent\textit{Editorial note.} Queue index 9 of queue\_defos.json (pb=3021, stmt 274-297, proof 299-373).\par\medskip
\noindent\textit{Allowed standard prerequisites (editorial).}\begin{itemize}
\item Explicit prerequisite boundary. The cited inputs are prerequisites; the argument in the retained span is not replaced by any cited result.
\end{itemize}\medskip
\ifdefined\hypertarget\hypertarget{sourcepointer-1}{}\fi
\noindent\textit{Source reference.} Original lemma 00S0: sources/stacks-defos/pinned-original-source.tex, source lines 373--373. Cited by the retained statement or proof but not part of the collected spans; supplied only as a navigation pointer. The argument does not depend on its content.\par\medskip
\ifdefined\hypertarget\hypertarget{sourcepointer-2}{}\fi
\noindent\textit{Source reference.} Original lemma 08S4: sources/stacks-defos/pinned-original-source.tex, source lines 260--267. Cited by the retained statement or proof but not part of the collected spans; supplied only as a navigation pointer. The argument does not depend on its content.\par\medskip
\begin{lemma}
\label{lemma-huge-diagram}
Given a commutative diagram
$$
\xymatrix{
& 0 \ar[r] & N_2 \ar[r] & B'_2 \ar[r] & B_2 \ar[r] & 0 \\
& 0 \ar[r]|\hole & I_2 \ar[u]_{c_2} \ar[r] &
A'_2 \ar[u] \ar[r]|\hole & A_2 \ar[u] \ar[r] & 0 \\
0 \ar[r] & N_1 \ar[ruu] \ar[r] & B'_1 \ar[r] & B_1 \ar[ruu] \ar[r] & 0 \\
0 \ar[r] & I_1 \ar[ruu]|\hole \ar[u]^{c_1} \ar[r] &
A'_1 \ar[ruu]|\hole \ar[u] \ar[r] & A_1 \ar[ruu]|\hole \ar[u] \ar[r] & 0
}
$$
with front and back solutions to (\ref{equation-to-solve}) we have
\begin{enumerate}
\item There exist a canonical element in
$\Ext^1_{B_1}(\NL_{B_1/A_1}, N_2)$
whose vanishing is a necessary and sufficient condition for the existence
of a ring map $B'_1 \to B'_2$ fitting into the diagram.
\item If there exists a map $B'_1 \to B'_2$ fitting into the diagram
the set of all such maps is a principal homogeneous space under
$\Hom_{B_1}(\Omega_{B_1/A_1}, N_2)$.
\end{enumerate}
\end{lemma}

\begin{proof}
Let $E = B_1$ viewed as a set.
Consider the surjection $A_1[E] \to B_1$ with kernel $J$ used
to define the naive cotangent complex by the formula
$$
\NL_{B_1/A_1} = (J/J^2 \to \Omega_{A_1[E]/A_1} \otimes_{A_1[E]} B_1)
$$
in
Algebra, Section \ref{algebra-section-netherlander}.
Since $\Omega_{A_1[E]/A_1} \otimes B_1$ is a free
$B_1$-module we have
$$
\Ext^1_{B_1}(\NL_{B_1/A_1}, N_2) =
\frac{\Hom_{B_1}(J/J^2, N_2)}
{\Hom_{B_1}(\Omega_{A_1[E]/A_1} \otimes B_1, N_2)}
$$
We will construct an obstruction in the module on the right.
Let $J' = \Ker(A'_1[E] \to B_1)$. Note that there is a surjection
$J' \to J$ whose kernel is $I_1A'_1[E]$.
For every $e \in E$ denote $x_e \in A_1[E]$ the corresponding variable.
Choose a lift $y_e \in B'_1$ of the image of $x_e$ in $B_1$ and
a lift $z_e \in B'_2$ of the image of $x_e$ in $B_2$.
These choices determine $A'_1$-algebra maps
$$
A'_1[E] \to B'_1 \quad\text{and}\quad A'_1[E] \to B'_2
$$
The first of these gives a map $J' \to N_1$, $f' \mapsto f'(y_e)$
and the second gives a map $J' \to N_2$, $f' \mapsto f'(z_e)$.
A calculation shows that these maps annihilate $(J')^2$.
Because the left square of the diagram (involving $c_1$ and $c_2$)
commutes we see that these maps agree on $I_1A'_1[E]$ as maps into $N_2$.
Observe that $B'_1$ is the pushout of $J' \to A'_1[E]$ and $J' \to N_1$. 
Thus, if the maps $J' \to N_1 \to N_2$ and $J' \to N_2$ agree, then we
obtain a map $B'_1 \to B'_2$ fitting into the diagram.
Thus we let the obstruction be the class of the map
$$
J/J^2 \to N_2,\quad f \mapsto f'(z_e) - \nu(f'(y_e))
$$
where $\nu : N_1 \to N_2$ is the given map and where $f' \in J'$
is a lift of $f$. This is well defined by our remarks above.
Note that we have the freedom
to modify our choices of $z_e$ into $z_e + \delta_{2, e}$
and $y_e$ into $y_e + \delta_{1, e}$ for some $\delta_{i, e} \in N_i$.
This will modify the map above into
$$
f \mapsto f'(z_e + \delta_{2, e}) - \nu(f'(y_e + \delta_{1, e})) =
f'(z_e) - \nu(f'(z_e)) +
\sum (\delta_{2, e} - \nu(\delta_{1, e}))\frac{\partial f}{\partial x_e}
$$
This means exactly that we are modifying the map $J/J^2 \to N_2$
by the composition $J/J^2 \to \Omega_{A_1[E]/A_1} \otimes B_1 \to N_2$
where the second map sends $\text{d}x_e$ to
$\delta_{2, e} - \nu(\delta_{1, e})$. Thus our obstruction is well defined
and is zero if and only if a lift exists.

\medskip\noindent
Part (2) comes from the observation that given two maps
$\varphi, \psi : B'_1 \to B'_2$ fitting into the diagram, then
$\varphi - \psi$ factors through a map $D : B_1 \to N_2$ which
is an $A_1$-derivation:
\begin{align*}
D(fg) & = \varphi(f'g') - \psi(f'g') \\
& =
\varphi(f')\varphi(g') - \psi(f')\psi(g') \\
& =
(\varphi(f') - \psi(f'))\varphi(g') + \psi(f')(\varphi(g') - \psi(g')) \\
& =
gD(f) + fD(g)
\end{align*}
Thus $D$ corresponds to a unique $B_1$-linear map
$\Omega_{B_1/A_1} \to N_2$. Conversely, given such a linear map
we get a derivation $D$ and given a ring map $\psi : B'_1 \to B'_2$
fitting into the diagram
the map $\psi + D$ is another ring map fitting into the diagram.
\end{proof}

\end{document}
